\section{Performance Analysis}

Beyond correctness, the correlator must meet performance requirements. This section quantifies throughput, memory, scalability, and resource utilization.

\subsection{Computational Throughput and Scalability}

For the baseline configuration (4 antennas, 256 channels), the correlator processes 10 seconds of data in 5.9 seconds, achieving 1.7$\times$ real-time performance. Smaller configurations (2 antennas, 64 channels) reach 4.8$\times$ real-time, while larger configurations (8 antennas or 1024 channels) fall below real-time (0.5–0.7$\times$).

Antenna scaling follows $O(N_{\text{ant}}^2)$ complexity due to baseline count $N(N-1)/2$. Processing time per baseline decreases with array size as fixed overhead (loading, F-Engine) amortises across more baselines.

Channel scaling shows $O(N_{\text{ch}} \log N_{\text{ch}})$ for FFT and linear for correlation. Time per channel plateaus at 18 ms for large $N$ where correlation dominates over FFT overhead.

\subsection{Memory Consumption and CPU Utilization}

Memory allocations include input buffer, channelized data, and visibility buffer. For 4 antennas, 256 channels: total $< 1$ GB. Scaling to 16 antennas, 1024 channels requires only a few gigabytes, feasible on commodity hardware. Memory scales linearly with antenna and channel count.

Profiling shows F-Engine FFT dominates (36\%), followed by X-Engine correlation (24\%). Together these constitute 60\% of processing, confirming DSP operations as the bottleneck. Data I/O comprises 19\%, suggesting optimization could help but is not limiting.

NumPy automatically leverages multi-core processors. Preliminary multiprocessing tests achieved 3.0$\times$ speedup on 4 cores (75\% efficiency) for F-Engine parallelization, limited by overhead. GPU acceleration could provide orders-of-magnitude improvement.
\subsection{Performance Summary}

The baseline configuration meets requirements: 1.7$\times$ real-time processing, $< 1$ GB memory, phase accuracy $< 0.2^\circ$, amplitude accuracy $< 1\%$. Scaling to 8 antennas falls slightly short of real-time but parallelization or GPU acceleration would address this.

Optimization opportunities include GPU acceleration (10–100$\times$ speedup potential), reduced precision (complex64 instead of complex128 halves bandwidth), streaming architecture (reduces memory), and polyphase filter banks (better resolution at similar cost).

The current performance suffices for research validation. Operational deployment for sustained observations would benefit from GPU acceleration and streaming data handling.
